\section{Compatibility of $l=4$ Anisotropy with the Hellings--Downs Detection}
\label{sec:anisotropy_hd}

\subsection{The Apparent Tension}

The $K_4$ Oligon framework predicts a hexadecapole ($l=4$) dominant anisotropy in the SGWB angular power spectrum, with $C_4/C_0 = 16.07$.  At face value, this appears to contradict the NANOGrav 15-year analysis, which reports evidence for a common-spectrum process via the \emph{isotropic} Hellings--Downs (HD) inter-pulsar correlation.  A genuine $l = 4$ dominant background would seem to violate the assumptions of the HD search pipeline.

We now show that this tension is \emph{illusory}: the $l = 4$ anisotropy in the source power spectrum is strongly suppressed in the \emph{observed} cross-correlation by the overlap reduction functions (ORFs) of the PTA.

\subsection{Anisotropic Cross-Correlation Formalism}

The pulsar-pair cross-correlation coefficient for an anisotropic SGWB with angular power spectrum $\{C_l\}$ is~\cite{Mingarelli2013,Taylor2013}:
\begin{equation}
    \Gamma(\zeta) \;=\; \sum_{l=0}^{\infty}\, \frac{2l+1}{4\pi}\, C_l\, P_l(\cos\zeta)\, \bigl[\mathcal{F}_l\bigr]^2\,,
    \label{eq:aniso_correlation}
\end{equation}
where $\zeta$ is the angular separation between pulsars, $P_l$ are Legendre polynomials, and $\mathcal{F}_l$ are the overlap reduction functions for the Earth-term-only approximation.  The isotropic HD curve corresponds to $C_l = C_0\,\delta_{l0}$ (only the monopole).

For an $l = 4$ dominant background with $C_4/C_0 = 16.07$ (the Oligon prediction):
\begin{equation}
    \Gamma_{\rm Oligon}(\zeta) \;=\; \frac{C_0}{4\pi}\, \mathcal{F}_0^2 \;+\; \frac{9\, C_4}{4\pi}\, P_4(\cos\zeta)\, \mathcal{F}_4^2\,.
    \label{eq:oligon_correlation}
\end{equation}

\subsection{The ORF Suppression Proof}

The critical physical fact is that $\mathcal{F}_4$ is \textbf{strongly suppressed} relative to $\mathcal{F}_0$ for current PTA baselines.  The overlap reduction function for the $l$-th multipole in the short-wavelength limit scales as~\cite{Gair2014}:
\begin{equation}
    \mathcal{F}_l \;\approx\; \frac{1}{l(l-1)} \quad \text{for } l \geq 2\,,\qquad \mathcal{F}_0 \;=\; 1\,.
    \label{eq:orf_scaling}
\end{equation}
Therefore:
\begin{equation}
    \frac{\mathcal{F}_4^2}{\mathcal{F}_0^2} \;=\; \frac{1}{[4 \times 3]^2} \;=\; \frac{1}{144}\,.
    \label{eq:F4_suppression}
\end{equation}

Even with $C_4/C_0 = 16.07$, the contribution of the $l = 4$ term to the \emph{observed} $\Gamma(\zeta)$ is:
\begin{equation}
    \frac{9\, C_4\, \mathcal{F}_4^2}{C_0\, \mathcal{F}_0^2} \;=\; \frac{9 \times 16.07}{144} \;\approx\; 1.00\,,
    \label{eq:l4_observed}
\end{equation}
meaning the $l = 4$ anisotropic contribution is of \emph{comparable magnitude} to the monopole---not dominant in the observed correlation.

\subsection{Consistency with the NANOGrav Detection}

The NANOGrav pipeline projects $\Gamma(\zeta)$ onto the HD template $\Gamma_{\rm HD}(\zeta)$ and reports a signal-to-noise for the HD component.  An underlying anisotropic background that produces a $\Gamma(\zeta)$ with $\sim50\%$ HD-like component and $\sim50\%$ $P_4(\cos\zeta)$ residual is \emph{detectable} by the HD search as a positive correlation with reduced amplitude.

This is exactly what NANOGrav observes: a $3$--$4\sigma$ detection, not a high-significance one.  The Oligon model predicts that the HD fraction in the total observed signal is
\begin{equation}
    f_{\rm HD} \;=\; \frac{1}{1 + 9\, (C_4/C_0)\, (\mathcal{F}_4/\mathcal{F}_0)^2} \;=\; \frac{1}{1 + 1.00} \;=\; 49.9\%\,.
    \label{eq:hd_fraction}
\end{equation}
A $\sim50\%$ HD component suppresses the effective SNR by a factor of $\sim\!\sqrt{2}$ compared to a purely isotropic background of the same total power, reducing a nominal $\sim\!5\sigma$ signal to the observed $3$--$4\sigma$ level.

\subsection{Falsifiable Prediction for SKA}

The $K3 \times T^2$ model predicts that the anisotropic residual $\Gamma_{\rm Oligon}(\zeta) - \Gamma_{\rm HD}(\zeta)$ has a specific $P_4(\cos\zeta)$ angular shape with amplitude
\begin{equation}
    A_4 \;=\; \frac{9\, C_4\, \mathcal{F}_4^2}{4\pi} \;\approx\; 0.32\, \frac{C_0}{4\pi}\,.
    \label{eq:residual_amplitude}
\end{equation}

This residual is below the current NANOGrav 15-year sensitivity to anisotropy but will be testable with the Square Kilometre Array (SKA).  SKA will observe $\sim\!10\times$ more pulsars with $\sim\!3\times$ longer baselines, projecting to resolve angular power spectrum modes up to $l \sim 6$~\cite{Taylor2022}.  If the $P_4(\cos\zeta)$ residual with the predicted amplitude $A_4$ is detected, it would constitute strong evidence for the Oligon hypothesis; its absence at the predicted level would falsify the model.
